% GENERATED FILE. Edit canonical lessons, then run npm run generate:books.
% canonical-source-hash: fnv1a64:1ea2d3a328387dfb
% canonical-lessons: FR-C337-admettre, FR-C337-consensus, FR-C337-compromis, FR-C337-ceder, FR-C337-conceder

\chapter{Trouver un accord}
\label{ch:fr-trouver-un-accord}
\hlchaptermodality{337}

\begin{chapteropening}
\textbf{By the end of this chapter:} \emph{I can admit someone is right, give way, concede a point and reach a consensus or a compromise.}

Complete the last lesson of chapter 337: I can admit someone is right, give way, concede a point and reach a consensus or a compromise.
\end{chapteropening}

\section[admettre]{admettre — to admit}

\label{lesson:FR-C337-admettre}



\begin{quote}
Before the new one: say the French for to retort, then the French for to contradict.
\end{quote}

\subsection*{You'll want to know: admettre}
\textbf{admettre} — \textquotedblleft{}to admit\textquotedblright{}.

\begin{quote}
Je dois admettre que ton idée était meilleure.
\end{quote}

\begin{quote}
\emph{I have to admit that your idea was better.}
\end{quote}

Conjugated like \emph{mettre}; past participle \emph{admis}. It also means to let in: \emph{être admis à l'université}.

\begin{cousinweb}
\emph{admettre} comes through Old French \emph{admettre} from Latin \emph{admittere} \textquotedblleft{}to let in, to allow\textquotedblright{} (EtymDB; Lewis \& Short), built from \emph{ad} \textquotedblleft{}to\textquotedblright{} and \emph{mittere} \textquotedblleft{}to send\textquotedblright{}: you let a fact in.
\end{cousinweb}

\begin{grammarlens}[title={What you've built}]
One more word for telling, explaining and describing at B1.
\end{grammarlens}

\subsection*{Your turn}
\begin{itemize}
\raggedright
  \item \emph{Say it:} \emph{admettre}
  \item \emph{Say it:} \emph{admettre}, once more
  \item \emph{Say it:} a sentence of your own with \emph{admettre}
  \item \emph{Recall:} say the French for to retort, then the French for to contradict, then say \emph{admettre} again
\end{itemize}

\begin{tcolorbox}[breakable,colback=teal!4,colframe=teal!35!black,title={Before you move on}]
What does admettre mean? (\textquotedblleft{}To admit\textquotedblright{}.) Say it once more.
\end{tcolorbox}
\section[le consensus]{le consensus — the consensus, broad agreement}

\label{lesson:FR-C337-consensus}



\begin{quote}
Before the new one: say the French for to contradict, then the French for to admit.
\end{quote}

\subsection*{You'll want to know: le consensus}
\textbf{le consensus} — \textquotedblleft{}the consensus, broad agreement\textquotedblright{}.

\begin{quote}
Après deux heures, nous avons trouvé un consensus.
\end{quote}

\begin{quote}
\emph{After two hours, we reached a consensus.}
\end{quote}

The final \emph{s} is said: \emph{consensuss}.

\begin{cousinweb}
\emph{consensus} was borrowed whole from Latin \emph{cōnsēnsus} \textquotedblleft{}agreement, accord\textquotedblright{} (EtymDB), from \emph{cōnsentīre} \textquotedblleft{}to agree\textquotedblright{}, built on \emph{sentīre} \textquotedblleft{}to feel\textquotedblright{} (Lewis \& Short).
\end{cousinweb}

\begin{grammarlens}[title={What you've built}]
One more word for telling, explaining and describing at B1.
\end{grammarlens}

\subsection*{Your turn}
\begin{itemize}
\raggedright
  \item \emph{Say it:} \emph{le consensus}
  \item \emph{Say it:} \emph{le consensus}, once more
  \item \emph{Say it:} a sentence of your own with \emph{le consensus}
  \item \emph{Recall:} say the French for to contradict, then the French for to admit, then say \emph{le consensus} again
\end{itemize}

\begin{tcolorbox}[breakable,colback=teal!4,colframe=teal!35!black,title={Before you move on}]
What does le consensus mean? (\textquotedblleft{}The consensus, broad agreement\textquotedblright{}.) Say it once more.
\end{tcolorbox}
\section[le compromis]{le compromis — the compromise}

\label{lesson:FR-C337-compromis}



\begin{quote}
Before the new one: say the French for to admit, then the French for the consensus.
\end{quote}

\subsection*{You'll want to know: le compromis}
\textbf{le compromis} — \textquotedblleft{}the compromise\textquotedblright{}.

\begin{quote}
Finalement, nous avons trouvé un compromis : un film, puis une pizza.
\end{quote}

\begin{quote}
\emph{In the end we found a compromise: a film, then a pizza.}
\end{quote}

\emph{Trouver un compromis} — to reach a compromise. The same form in the plural.

\begin{cousinweb}
\emph{compromis} comes from Latin \emph{compromissum} \textquotedblleft{}a mutual promise to abide by an arbiter's award\textquotedblright{} (Lewis \& Short; EtymDB), built on \emph{prōmittere} \textquotedblleft{}to promise\textquotedblright{}.
\end{cousinweb}

\begin{grammarlens}[title={What you've built}]
One more word for telling, explaining and describing at B1.
\end{grammarlens}

\subsection*{Your turn}
\begin{itemize}
\raggedright
  \item \emph{Say it:} \emph{le compromis}
  \item \emph{Say it:} \emph{le compromis}, once more
  \item \emph{Say it:} a sentence of your own with \emph{le compromis}
  \item \emph{Recall:} say the French for to admit, then the French for the consensus, then say \emph{le compromis} again
\end{itemize}

\begin{tcolorbox}[breakable,colback=teal!4,colframe=teal!35!black,title={Before you move on}]
What does le compromis mean? (\textquotedblleft{}The compromise\textquotedblright{}.) Say it once more.
\end{tcolorbox}
\section[céder]{céder — to give in, to give way; to hand over}

\label{lesson:FR-C337-ceder}



\begin{quote}
Before the new one: say the French for the consensus, then the French for the compromise.
\end{quote}

\subsection*{You'll want to know: céder}
\textbf{céder} — \textquotedblleft{}to give in, to give way; to hand over\textquotedblright{}.

\begin{quote}
Il a insisté, et j'ai fini par céder.
\end{quote}

\begin{quote}
\emph{He insisted, and I ended up giving in.}
\end{quote}

\emph{Céder à quelqu'un} — to give in to someone. \emph{Je cède}, with \emph{è} before a silent ending.

\begin{cousinweb}
\emph{céder} is a learned borrowing from Latin \emph{cēdere} \textquotedblleft{}to go, to withdraw, to yield\textquotedblright{} (EtymDB; Lewis \& Short).
\end{cousinweb}

\begin{grammarlens}[title={What you've built}]
One more word for telling, explaining and describing at B1.
\end{grammarlens}

\subsection*{Your turn}
\begin{itemize}
\raggedright
  \item \emph{Say it:} \emph{céder}
  \item \emph{Say it:} \emph{céder}, once more
  \item \emph{Say it:} a sentence of your own with \emph{céder}
  \item \emph{Recall:} say the French for the consensus, then the French for the compromise, then say \emph{céder} again
\end{itemize}

\begin{tcolorbox}[breakable,colback=teal!4,colframe=teal!35!black,title={Before you move on}]
What does céder mean? (\textquotedblleft{}To give in, to give way; to hand over\textquotedblright{}.) Say it once more.
\end{tcolorbox}
\section[concéder]{concéder — to concede, to grant}

\label{lesson:FR-C337-conceder}



\begin{quote}
Before the new one: say the French for the compromise, then the French for to give in.
\end{quote}

\subsection*{You'll want to know: concéder}
\textbf{concéder} — \textquotedblleft{}to concede, to grant\textquotedblright{}.

\begin{quote}
Je te concède ce point : tu avais raison.
\end{quote}

\begin{quote}
\emph{I grant you that point: you were right.}
\end{quote}

\emph{Je concède}, like \emph{céder}. Formal; in speech, \emph{je reconnais que} is common.

\begin{cousinweb}
\emph{concéder} comes from Latin \emph{concēdere} \textquotedblleft{}to withdraw, to grant\textquotedblright{} (EtymDB; Lewis \& Short), a strengthened form of \emph{cēdere} \textquotedblleft{}to go, to yield\textquotedblright{}.
\end{cousinweb}

\begin{grammarlens}[title={What you've built}]
One more word for telling, explaining and describing at B1.
\end{grammarlens}

\subsection*{Your turn}
\begin{itemize}
\raggedright
  \item \emph{Say it:} \emph{concéder}
  \item \emph{Say it:} \emph{concéder}, once more
  \item \emph{Say it:} a sentence of your own with \emph{concéder}
  \item \emph{Recall:} say the French for the compromise, then the French for to give in, then say \emph{concéder} again
\end{itemize}

\begin{tcolorbox}[breakable,colback=teal!4,colframe=teal!35!black,title={Before you move on}]
What does concéder mean? (\textquotedblleft{}To concede, to grant\textquotedblright{}.) Say it once more.
\end{tcolorbox}
